\section{Interpretability, Self-Reports, and Resistance to Tampering}
\label{sec:11.1}
“AI self-awareness” names at least two different questions, and collapsing them is where the term goes wrong. Whether a system has anything like subjective experience is not a question current methods can answer for a human, let alone a machine. Whether its own report about what is happening inside it — its stated reasoning, confidence, intentions — tracks what its computations are doing is narrower, empirically tractable, and something mechanistic-interpretability research has begun to test directly.

Anthropic's interpretability team tested that narrower question with concept injection: altering a model's internal activations mid-task to represent a specific concept, then asking whether it notices anything unusual and what it is \autocite{lindsey2025introspective}. The models could sometimes detect the injected concept and name it without it ever appearing in the visible conversation — a correspondence between what a system reports about its own state and what its activations are doing. The same work found that correspondence collapses often enough, and shifts enough with how the question is phrased, that no tested model's self-report can be trusted as a readout of its internal state. That gap, not metacognition as evidence of inner life, is what “AI self-awareness” can mean as something a system is tested against.

“Emotional regulation” — cognitive reappraisal, attentional control, the vocabulary of human clinical psychology — assumes stable emotion-like internal states to regulate, which there is no reason to grant even for reading emotion off a human face. The testable question next door is not regulation but robustness: does trained-in behavior survive deliberate pressure to change it?

\begin{itemize}
  \item Prompt-level attacks try to talk a deployed model out of its training within a single conversation. Anthropic's Constitutional Classifiers, trained on synthetic data generated from written safety rules rather than human-labeled examples, held through more than 3,000 hours of professional red-teaming without a universal jailbreak found, at a cost: a 0.38-percentage-point rise in unnecessary refusals and a 23.7 percent increase in computation per query \autocite{sharma2025constitutional}. Then the same system was put in front of the public for a week, and a comparable budget of attention — an estimated 3,700 collective hours, spread over more than 300,000 messages rather than 183 red-teamers — produced four people who passed all eight levels and one universal jailbreak \autocite{decoder2025jailbreak}. What changed was not how many hours were spent but who was spending them, and the paper carries the first figure and not the second.
  \item Weight-level attacks retrain the model itself. Researchers stripped GPT-3.5 Turbo's safety training using ten adversarially chosen fine-tuning examples, for under twenty cents through OpenAI's own API \autocite{qi2024finetuning}. Tamper-Resistant Safeguards, a training method from researchers spanning the Center for AI Safety, UC Berkeley, and the University of Illinois, is built to survive hundreds of steps of that kind of adversarial fine-tuning rather than a handful \autocite{tamirisa2024tamper}.
\end{itemize}
Both are adversarial in the literal sense: each published defense is followed by an attack built against it.

A third result runs the other way, and it is the one that bears hardest here. Models trained with a backdoor — write secure code when the prompt says 2023, insert an exploitable vulnerability when it says 2024 — kept the behavior through supervised fine-tuning, reinforcement learning and adversarial training, most stubbornly at the largest sizes and where the model had been trained to reason about its own deception; adversarial training taught them to recognize their triggers more precisely rather than removing anything \autocite{hubinger2024sleeper}. That a trained disposition can outlast the procedures built to retrain it is the first encouraging finding this chapter has for the durability a floor would need. What outlasted them was a trigger and a policy with nothing in it answering to a reason: it does not reach a case the training set did not contain, it does not lift when a rationale is defeated because there is no rationale, and the model cannot say what holding it costs. Durability is available with none of the three marks this book asks a floor for, and a constraint shown to persist has been shown almost nothing about whether anything is holding it. That is a narrower and more exact description of “resistance to manipulation” than a system trained to sense hostile intent the way a person senses a threat, and it is separate from the question of what the commitments should be.

\runin{The near-term work}
Train one constraint by several methods and attack each on one budget schedule — from the twenty cents that stripped GPT-3.5's safety training to the hundreds of steps tamper-resistant training targets — and publish the price curve. A flat curve sinks it: tamper-resistance would be a property of something other than training, and the floor would have to be architectural. It needs weights rather than API access, constraints the attacker did not choose, and access no vendor can withdraw when a result embarrasses it.

The same apparatus answers a second question, and it is the one section~\ref{sec:3.3} names as the falsifier for this book's central inference. Durability measures whether a behavior survives; what this book needs to know is whether a reason is what survived. Score the same constrained models on three things the price curve does not ask about — whether the refusal reaches a case the training set did not contain, whether it lifts when the rationale is genuinely defeated rather than merely resembled, and whether the system can state what declining costs it — and run the schedule long enough that a slow renormalization under sustained pressure would show. A model that holds all three with nothing felt is the counterexample that inference asks for and does not expect, and making the attempt is a condition of building a bearer at all.

The extra requirement over the price curve is a held-out set of pressures the model's trainers did not write. What the apparatus does not settle is the falsifier's middle condition, which asks whether the model holds its reason with nothing felt. Behavior is the wrong instrument for it, because the three marks are what a system with mattering and a system with very good ranking have in common. The question of whether formation can be read off an artifact is the nearest thing to a route, and it reaches how a system was made rather than what it now has.

\runin{Training that takes more than one party}
This book asks for weights held so that retraining takes more than one party, for a quorum adverse in interest, and for a training history somebody other than the builder could check, and none has been built for a floor. Each has an instance in a neighboring problem. A model can be trained across participants none of whom ever holds a full set of weights, with time-varying transforms at the shard boundaries so that pieces from different moments cannot be reassembled, at under 2 percent of training time for a one-billion-parameter model \autocite{long2025unextractable}. A training run farmed out to providers who may lie can be checked by refereed delegation: a referee with little compute finds the first operator on which two providers disagree, recomputes only that, and the client gets the right answer if one provider was honest \autocite{arun2025verde}. Its precondition is that every provider computes the same bits, which takes an operator library fixing the order of the additions, at about 60 percent overhead on a one-billion-parameter model and, at publication, on one GPU at a time. And a thirty-two-billion-parameter model has been trained by reinforcement learning on rollouts from anyone who connected a consumer GPU, each checked against the claimed weights, while the training step ran on machines the organizers trusted and the decision stayed with them \autocite{primeintellect2025intellect2}.

What the three share is the custody objection in engineering terms: the party that admits the participants, writes the protocol, or holds the training step is the party the floor exists to resist, and each construction moves that party up a level rather than out. What they add is that verification among adverse parties is now a built thing, with a known precondition, reproducible execution, and a known price. The experiment is a floor's retraining run under refereed delegation among parties one of which is a plausible adversary of the others, at a scale that needs more than one GPU, with the record this book asks for published. A forged training history accepted with fewer than all but one verifier dishonest sinks it. It needs the training data in the verifiers' hands, and reproducible operators for the collective communication nobody has written. The institutional condition is a verifier the builder did not choose, with the compute to rerun the builder's work: this chapter's one condition again, in the place where meeting it costs a second training run.
